\documentclass[11pt]{article}
\usepackage[margin=1in]{geometry}
\usepackage{amsmath,amssymb}
\usepackage{booktabs}
\usepackage{graphicx}
\usepackage[hidelinks]{hyperref}
\usepackage{times}

\title{\textbf{Sensitivity Analysis Protocol}\\[4pt]
\large Validation Methodology for the Kuramoto Phase Engine in
NEUROCARDIA\\[2pt]
\normalsize Companion document to the NEUROCARDIA White Paper\\[2pt]
\small Version 0.2 --- Protocol and Results}
\author{Americo Simoes \\ \texttt{amexsimoes@gmail.com}}
\date{October 14, 2026}

\begin{document}
\maketitle

\begin{center}
\small
Copyright \textcopyright{} 2026 Americo Simoes. All rights reserved.\\
This document is released under the Creative Commons Attribution 4.0
International License (CC-BY-4.0). The software it describes is
released under the GNU General Public License, version 3 or later
(GPL-3.0-or-later).
\end{center}

\vspace{1em}
\noindent\fbox{\parbox{\dimexpr\linewidth-2\fboxsep-2\fboxrule}{%
\textbf{Revision history.} This is revision 0.2 of the
\emph{Sensitivity Analysis Protocol}. Revision 0.1 (October 14, 2026)
specified the protocol, the five experimental arms, and the acceptance
criteria, but reserved the results section for a subsequent revision.
Revision 0.2 populates the results section, adds the residual AUC as a
fourth evaluation metric, introduces a companion acceptance criterion
for the residual diagnostic, and reports the finding that trajectory
robustness and residual robustness are disjoint. The protocol
specification of revision 0.1 is unchanged.}}

\vspace{1em}

\begin{abstract}
The NeuroCardia diagnostic engine applies a Kuramoto network of $N = 16$
coupled oscillators to the RR tachogram of a single-lead ECG recording,
producing a Cardiac Synchronization Index (CSI, $r(t)$) and a family of
early warning signals (EWS). Three parameter choices distinguish this
engine from the standard Kuramoto literature: a deterministic frequency
distribution drawn from the non-trivial Riemann zeros, a carrier
frequency five orders of magnitude above the cardiac band, and
logarithmic rather than linear spacing between adjacent frequencies.
This document specifies the formal sensitivity analysis protocol that
determines whether the resulting CSI and EWS are robust against these
choices, and reports the results of running the protocol against the
MIT-BIH Arrhythmia Database and a paired corpus of normal sinus and
atrial fibrillation recordings.

The principal findings are two. First, the absolute value of the CSI is
not robust to the frequency distribution at any coupling strength where
the metric has dynamic range. The mean Pearson correlation between the
CSI trajectory under the Riemann-zero configuration and under any
alternative configuration is at most $0.17$ at $K = 2$, and reaches the
pre-registered threshold of $0.80$ only for the linear-spacing
distribution at $K = 30$, where the network is saturated. Second, the
residual diagnostic that the companion papers report is not robust to
the frequency distribution either: three of five distributions exhibit
the residual at $K = 2$ (AUC 0.84--0.89), and two do not (AUC
0.54--0.56). The two robustness properties are disjoint: no monotone
relation exists between trajectory correlation and residual AUC across
the five arms.

The protocol's purpose is both scientific defensibility and
satisfaction of the analytical validation requirements of the
IEC 62304 / FDA 510(k) pathway.
\end{abstract}

\tableofcontents

\section{Introduction}

\subsection{Context}

The NeuroCardia diagnostic engine applies a Kuramoto model of coupled
phase oscillators to the RR interval series of a single-lead ECG
recording. The model treats the sinoatrial node as a population of
coupled pacemaker cells and produces a time-varying order parameter
$r(t)$, called the Cardiac Synchronization Index (CSI). A family of
derived metrics --- variance, lag-1 autocorrelation, skewness,
kurtosis --- is computed over sliding windows and reported as early
warning signals (EWS) of impending rhythm transitions.

The engine is a specific parameterization of the standard Kuramoto
model. Three of its parameter choices are inherited from the
Convergent reference framework and are not found in the standard
literature. This document specifies the protocol by which those
choices are tested for robustness, and reports the results.

\subsection{Why sensitivity analysis}

A diagnostic metric is only as trustworthy as its robustness to
arbitrary choices. If the CSI is driven by the physical dynamics of the
cardiac rhythm, it should be stable across reasonable variations of the
model's free parameters. If it is driven by the specific values of
those parameters, then any diagnostic claim derived from it is
conditioned on those values, and the model's scientific value is
correspondingly reduced.

For a research tool, this is an interesting question. For a medical
device, it is a regulatory requirement. The FDA's guidance on
analytical validation for SaMD requires that the manufacturer
demonstrate that the algorithm computes what it claims to compute, and
that its outputs are not artifacts of parameter selection.

\subsection{Scope}

This document specifies: the three parameter choices under test
(Section~2), the experimental design comprising five arms
(Section~3), the evaluation metrics applied to each arm (Section~4),
the acceptance and falsification criteria (Section~5), the reporting
protocol (Section~6), and the results of running the protocol
(Section~7).

\section{Parameters Under Test}

\subsection{Deterministic frequency distribution}

Standard Kuramoto implementations assign each oscillator a natural
frequency drawn from a probability distribution $g(\omega)$ ---
commonly Gaussian, uniform, or Lorentzian. NeuroCardia instead assigns
frequencies deterministically from the non-trivial zeros of the Riemann
zeta function. The first 16 such zeros are used for the production
network of $N = 16$ oscillators, and the first 24 are tabulated for
extended networks.

The author does not claim that the Riemann zeros have a cardiac
physiological correlate. The choice reflects a mathematical rather
than a physiological motivation. Whether this matters for the
diagnostic output is the question this protocol addresses.

\subsection{Carrier frequency scale}

The reference parameterization places the carrier frequency at
$\Omega_T = 587{,}000$ rad/s, five orders of magnitude above the cardiac
rhythm (approximately $1$ Hz). The dimensionless order parameter $r(t)$
is, by construction, invariant to the absolute scale of $\Omega_T$,
because the RR-to-phase mapping and the R-peak phase resets couple the
network to relative beat timing only. The normalized frequencies used in
production are obtained by subtracting the carrier and scaling to unit
variance; this linear transformation does not change the shape of the
distribution.

\subsection{Logarithmic frequency spacing}

The Riemann zeros grow denser as $\gamma$ increases. Their spacing is
governed by the Riemann-von Mangoldt formula, giving an approximately
logarithmic density. This contrasts with the linear spacing of a
uniformly distributed frequency set, and with the unimodal clustering
of a Gaussian set. Whether the CSI depends on the shape of the
frequency distribution, or only on its overall variance, is an
empirical question.

\section{Experimental Design}

\subsection{Benchmark corpus}

The analysis runs on a curated subset of the PhysioNet reference
databases chosen to span the clinically relevant range of cardiac
dynamics. For the trajectory sensitivity analysis, 48 records from the
MIT-BIH Arrhythmia Database~\cite{mitdb} were used. For the residual
sensitivity analysis, a paired corpus of nine records from the MIT-BIH
Normal Sinus Rhythm Database~\cite{nsrdb} and nine records from the
MIT-BIH Atrial Fibrillation Database~\cite{afdb} was used, each trimmed
to the first ten minutes.

All records are annotated. R-peak detection is performed by the
Pan-Tompkins module in NeuroCardia~\cite{pan1985}, then verified against
the reference annotations before any CSI computation. Records with
$Se < 99\%$ or $PPV < 99\%$ are excluded from the analysis and reported
separately.

\subsection{Experimental arms}

Five arms share an identical preprocessing pipeline (Pan-Tompkins
R-peak detection, RR-to-phase mapping) and differ only in the Kuramoto
engine frequency distribution.

\begin{table}[h]
\centering
\begin{tabular}{lll}
\toprule
Arm & Distribution & Notes \\
\midrule
A & Riemann zeros & Reference configuration \\
B & Uniform random & Stochastic; seed 42 \\
C & Gaussian & Deterministic (fixed seed) \\
D & Physiological scale & Bounded, uniform on a physiological range \\
E & Linear spacing & Uniform on an integer grid \\
\bottomrule
\end{tabular}
\caption{The five experimental arms. All arms use $N = 16$,
all-to-all coupling, and the same RR-to-phase mapping.}
\label{tab:arms}
\end{table}

\subsection{Controlled variables}

The number of oscillators ($N = 16$), the coupling topology
(all-to-all), the integration scheme (Euler with $\Delta t = 0.05$ s),
the phase-reset rule (reset oscillator 0 at each R-peak to the
cumulative phase), and the R-peak detector (Pan-Tompkins) are held
fixed across arms.

\section{Evaluation Metrics}

\subsection{Trajectory correlation}

The primary metric is the Pearson correlation between the CSI
trajectory $r_A(t)$ under Arm A and the trajectory $r_X(t)$ under each
alternative arm $X$. Because the trajectories may have different
lengths if the reset schedule differs, we sample both at the same beat
indices and compute the correlation over the shorter of the two.

The secondary trajectory metric is the normalized dynamic time warping
(DTW) distance between the two trajectories. Lower values indicate
greater similarity.

\subsection{Residual AUC}

The trajectory correlation measures whether the CSI is well-defined
under parameter perturbation. It does not measure whether the CSI is
diagnostic. The companion papers~\cite{simoes2026residual} identify a
specific diagnostic --- the residual of the CSI after regressing out
the coefficient of variation of the RR series --- and report that it
separates sinus from AF at $K = 2$ with AUC $0.89$ under Arm A.

We therefore add a second criterion: the residual AUC must be $\geq
0.80$ for the arm to be considered diagnostically usable. This
criterion is independent of the trajectory criterion, and
Section~7.4 shows that the two are disjoint.

\subsection{Bifurcation and transition timing}

Where the downstream diagnostic output is a discrete classification
over a sliding window, the concordance between Arm A and each
alternative arm is measured by Cohen's Kappa. This metric was specified
in revision 0.1 as a secondary criterion. It is reported where
applicable; the pilot corpus does not have the annotated transitions
required for a full analysis.

\section{Acceptance and Falsification Criteria}

\subsection{Trajectory acceptance}

The engine is deemed robust to a parameter if, across all records in the
benchmark corpus, the trajectory correlation with Arm A is at least
$r = 0.80$, and the normalized DTW distance is at most the within-Arm-A
variability measured on repeated runs of the same record.

\subsection{Residual acceptance}

The engine is deemed diagnostically robust to a parameter if, on the
paired sinus/AF corpus, the residual AUC for the corresponding arm is
at least $0.80$ at the optimum $K$.

\subsection{Partial robustness}

If the trajectory correlation falls between $r = 0.60$ and $r = 0.80$,
or the residual AUC falls between $0.65$ and $0.80$, the parameter is
deemed partially robust. In this case the white paper is updated to
qualify the corresponding claim, and the parameter is flagged as a
modeling assumption rather than a physical invariant.

\subsection{Falsification}

If the trajectory correlation falls below $r = 0.60$, or the residual
AUC falls below $0.65$, the parameter is deemed not robust for the
corresponding criterion. In this case: the white paper's corresponding
claim is restricted; the parameterization is designated a
model-dependent projection; and the diagnostic value of the metric is
reported as conditioned on the specific parameterization.

\subsection{Interpretation}

The protocol is designed to be falsifiable. A result that confirms
robustness is a positive scientific finding. A result that refutes
robustness is not a failure of the project --- it is an honest
characterization of what the metric measures and what it does not. The
white paper's regulatory claims are conditional on this outcome.

\section{Reporting Protocol}

\subsection{Pre-registration}

The protocol specified in this document is fixed before any sensitivity
analysis is run. Subsequent modifications are logged in a revision
history, dated, and accompanied by a stated rationale. The purpose is
to prevent post-hoc adjustment of the acceptance criteria in response
to observed results.

\subsection{Outputs}

For each record, the analysis produces: a CSV file with one row per arm
per seed per metric; a per-record CSV for the residual with one row per
arm per record; and a summary table aggregating across the corpus. All
outputs are committed to the repository under
\texttt{validation/reports/} and are reproducible with a single Gradle
command.

\subsection{Publication}

The results of the sensitivity analysis are reported in Section~7 of
this document. Where the results are of general scientific interest
they are additionally reported in the companion
documents~\cite{simoes2026residual,simoes2026engine}.

\section{Results}

\subsection{Trajectory robustness}

Table~\ref{tab:trajectory} reports the mean $\pm$ SD of the Pearson
correlation between the CSI trajectory under Arm A and under each
alternative arm, at five values of the coupling strength $K$, across
48 mitdb records.

\begin{table}[h]
\centering
\begin{tabular}{r|cccc}
\toprule
$K$ & Arm B & Arm C & Arm D & Arm E \\
\midrule
2.0 & $0.052 \pm 0.030$ & $0.076 \pm 0.042$ & $0.051 \pm 0.035$ & $0.168 \pm 0.089$ \\
5.0 & $0.250 \pm 0.045$ & $0.230 \pm 0.031$ & $0.250 \pm 0.045$ & $0.611 \pm 0.177$ \\
10.0 & $0.386 \pm 0.102$ & $0.366 \pm 0.085$ & $0.386 \pm 0.102$ & $0.778 \pm 0.171$ \\
30.0 & $0.458 \pm 0.165$ & $0.428 \pm 0.152$ & $0.458 \pm 0.165$ & $\mathbf{0.869 \pm 0.131}$ \\
100.0 & $-0.000 \pm 0.005$ & $0.000 \pm 0.006$ & $0.001 \pm 0.005$ & $0.001 \pm 0.005$ \\
\bottomrule
\end{tabular}
\caption{Mean $\pm$ SD Pearson correlation between the CSI trajectory
under Arm A and under each alternative arm, across 48 mitdb records.
The pre-registered threshold is $r = 0.80$. Only Arm E at $K = 30$
meets it, and that is the coupling strength at which the network is
saturated.}
\label{tab:trajectory}
\end{table}

Three observations.

First, at $K = 2$ --- the coupling strength at which the residual
diagnostic is strongest --- all four alternative arms fail the
trajectory criterion, with mean correlations of $0.05$--$0.17$ and
standard deviations small relative to the means. Every individual
record fails the threshold, not just the mean.

Second, the trajectory correlation rises monotonically with $K$ up to
$K = 30$, at which point Arm E meets the $0.80$ threshold. But $K =
30$ is the coupling strength at which the CSI saturates: the mean
order parameter across the corpus is $\approx 0.998$ with variance
$\approx 10^{-5}$, and the metric is a monotone function of RR
variability. Meeting the trajectory criterion at a saturated coupling
strength is not a robustness success.

Third, at $K = 100$ the trajectory correlation collapses to zero and
the DTW distance is the largest in the table. This is the regime in
which the Euler integration aliases the sine term: the phase advances
by more than $\pi$ per step and the sine evaluation is effectively
random.

The conclusion is that the CSI's absolute trajectory is not robust to
the frequency distribution at any coupling strength where the metric
has dynamic range.

\subsection{Residual robustness}

Table~\ref{tab:residual-arms} reports the residual AUC at $K = 2$ for
each of the five arms, on the paired sinus/AF corpus of nine and nine
records.

\begin{table}[h]
\centering
\begin{tabular}{llrrr}
\toprule
Arm & Distribution & Trajectory $\rho$ vs A & Residual AUC & Residual $p$ \\
\midrule
A & Riemann zeros & $1.000$ (ref) & $\mathbf{0.889}$ & $0.002$ \\
B & Uniform random & $0.052$ & $0.543$ & $0.986$ \\
C & Gaussian & $0.076$ & $\mathbf{0.852}$ & $0.007$ \\
D & Physiological scale & $0.051$ & $0.556$ & $0.982$ \\
E & Linear spacing & $0.168$ & $\mathbf{0.840}$ & $0.010$ \\
\bottomrule
\end{tabular}
\caption{Trajectory correlation and residual AUC at $K = 2$ for each
arm. Arms A, C, and E exhibit a diagnostic residual; arms B and D do
not. The trajectory correlation and residual AUC are not monotonically
related across arms.}
\label{tab:residual-arms}
\end{table}

Three of five arms exhibit the residual; two do not. The two
distributions that fail (uniform random and physiological scale) are
both bounded-support distributions, but the linear-spacing arm is also
bounded-support and does exhibit the residual, so bounded support is
not the distinguishing feature. We do not offer a mechanistic
explanation for the split.

The practical consequence is that the frequency distribution is part of
the metric's definition, and its choice matters for both the trajectory
and the residual.

\subsection{Orthogonality of the residual to standard HRV}

To test whether the residual carries information beyond standard HRV,
we regress it on the two features it is most correlated with in the
pilot corpus --- the Poincar\'e SD1/SD2 ratio and the coefficient of
variation of the RR series --- and compute the AUC of the
orthogonalized residual.

\begin{table}[h]
\centering
\begin{tabular}{lrr}
\toprule
Quantity & $\rho$ vs label & AUC \\
\midrule
Raw residual & $+0.675$ & $0.889$ \\
Residual $\perp$ $(\mathrm{sd1/sd2}, \mathrm{cv})$ & $+0.375$ & $0.716$ \\
\bottomrule
\end{tabular}
\caption{Orthogonality of the residual to standard HRV features. The
residual retains an AUC of $0.716$ after orthogonalization, down from
$0.889$.}
\label{tab:orthogonality}
\end{table}

The orthogonalized residual retains AUC $0.716$. Approximately two
thirds of the residual's diagnostic value is attributable to its linear
relationship with the two HRV features; one third is not. At $n = 18$
the confidence interval on the orthogonalized AUC is wide, and the
difference between $0.889$ and $0.716$ is not statistically significant
at this sample size.

\subsection{Window-length dependence of the residual}

The residual criterion of Section~4.2 was applied to full-record
analyses. A subsequent window-length sweep shows that the residual is a
record-level quantity: at 30-second windows the AUC is 0.50; at 1--5
minute windows it plateaus at 0.66--0.68; at the full record it is
0.83--0.89. The residual's within-record variance is 1.2--2.5 times
its between-record variance, and its window-to-window variation does
not correlate with window-level P-wave features. The criterion
therefore validates a record-level diagnostic, not a within-record
state tracker. The scope of the acceptance criterion is noted here for
clarity; it does not change the per-arm results reported below.

\subsection{Disjointness of trajectory and residual robustness}

Table~\ref{tab:disjoint} reproduces the key columns of
Tables~\ref{tab:trajectory} and~\ref{tab:residual-arms} side by side.

\begin{table}[h]
\centering
\begin{tabular}{lrr}
\toprule
Arm & Trajectory $\rho$ vs A (K = 2) & Residual AUC (K = 2) \\
\midrule
A & $1.000$ & $0.889$ \\
B & $0.052$ & $0.543$ \\
C & $0.076$ & $0.852$ \\
D & $0.051$ & $0.556$ \\
E & $0.168$ & $0.840$ \\
\bottomrule
\end{tabular}
\caption{Trajectory correlation and residual AUC are disjoint. Arm C
has the second-lowest trajectory correlation and the second-highest
residual AUC. Arm E has the highest trajectory correlation among the
alternative arms and the second-highest residual AUC. Arms B and D
fail on both metrics. No monotone relation exists between the two.}
\label{tab:disjoint}
\end{table}

The pre-registered trajectory criterion, applied to the useful
coupling regime, would report a failure for every arm. The residual
criterion reports a failure for two of five. The two criteria do not
identify the same arms as robust, and neither criterion subsumes the
other.

\section{Discussion}

\subsection{What the protocol establishes}

The protocol establishes two facts. First, the absolute value of the
CSI is not robust to the frequency distribution at any useful coupling
strength. This is a negative result about the metric's definition.
Second, the residual diagnostic that the companion papers report is
also not distribution-agnostic: three of five tested distributions
exhibit it.

Both facts condition the metric's interpretation. The CSI is a
parameterized quantity; its value refers to a specific
parameterization.

\subsection{What the protocol does not establish}

The protocol does not establish which distributions are physiologically
appropriate, nor whether any of the tested distributions has a claim to
physical significance. It does not establish the residual's behavior
under nonlinear regression, under alternative orthogonality procedures,
or on the full reference corpora. These are subjects for further work.

\subsection{Implications for the regulatory pathway}

For the FDA 510(k) pathway, the protocol demonstrates that the
algorithm's absolute output depends on parameter choices in a
documented way, and that the parameter dependence is quantified. The
regulatory strategy described in the white paper treats the CSI as a
research feature whose interpretation is conditioned on the
parameterization, and the residual as a separately validated
discriminative feature. Both postures are consistent with the results
reported here.

\section{Conclusion}

This document specifies and reports a formal sensitivity analysis of
the NeuroCardia Kuramoto phase engine against the PhysioNet reference
databases. The absolute value of the CSI is not robust to the frequency
distribution at any coupling strength where the metric has dynamic
range. The residual diagnostic at $K = 2$ is exhibited by three of five
tested frequency distributions. Trajectory robustness and residual
robustness are disjoint properties of the metric, and the protocol's
pre-registered trajectory criterion is not sufficient to validate the
diagnostic.

The results support the white paper's regulatory posture: the CSI is a
research metric whose value is parameterized; the residual is a
candidate diagnostic whose robustness is narrower than initially
reported but is quantifiable and reproducible.

\begin{thebibliography}{9}

\bibitem{mitdb}
G.~B.~Moody and R.~G.~Mark. The impact of the MIT-BIH Arrhythmia
Database. \emph{IEEE Engineering in Medicine and Biology Magazine},
20(3):45--50, 2001.

\bibitem{nsrdb}
R.~G.~Mark and G.~B.~Moody. MIT-BIH Normal Sinus Rhythm Database.
PhysioNet, 1990. \url{https://physionet.org/content/nsrdb/}.

\bibitem{afdb}
G.~B.~Moody and R.~G.~Mark. MIT-BIH Atrial Fibrillation Database.
PhysioNet, 1983. \url{https://physionet.org/content/afdb/}.

\bibitem{pan1985}
J.~Pan and W.~J.~Tompkins. A real-time QRS detection algorithm.
\emph{IEEE Transactions on Biomedical Engineering},
BME-32(3):230--236, 1985.

\bibitem{kuramoto1975}
Y.~Kuramoto. Self-entrainment of a population of coupled non-linear
oscillators. In \emph{International Symposium on Mathematical Problems
in Theoretical Physics}, pages 420--422. Springer, 1975.

\bibitem{acebron2005}
J.~A.~Acebr\'on, L.~L.~Bonilla, C.~J.~P\'erez Vicente, F.~Ritort, and
R.~Spigler. The Kuramoto model: A simple paradigm for synchronization
phenomena. \emph{Reviews of Modern Physics}, 77(1):137--185, 2005.

\bibitem{scheffer2009}
M.~Scheffer et al. Early-warning signals for critical transitions.
\emph{Nature}, 461(7260):53--59, 2009.

\bibitem{simoes2026engine}
A.~Simoes. The Kuramoto Engine: Physics and Mathematics of the Cardiac
Synchronization Index. Technical reference, revision 0.2, 2026.
Companion document.

\bibitem{simoes2026residual}
A.~Simoes. The Kuramoto Residual: Phase-Domain Information in the
Cardiac Rhythm Beyond Heart Rate Variability. Research paper, revision
0.2, 2026. Companion document.

\bibitem{simoes2026whitepaper}
A.~Simoes. NeuroCardia: An Open-Source, On-Device Phase-Domain
Diagnostic Engine for Single-Lead ECG. White paper, revision 0.2,
2026. Companion document.

\end{thebibliography}

\end{document}
